\documentclass[10pt,a4paper]{amsart}
\usepackage[margin=19mm]{geometry}
\usepackage{amsmath,amssymb,mathtools}
\usepackage[scr=rsfs,cal=euler]{mathalfa}
\usepackage{xfrac}
\usepackage[unicode,hidelinks,bookmarks=false]{hyperref}
\newcommand{\ie}{i.e.\ }
\newcommand{\cf}{cf.\ }
\newcommand{\resp}{respectively\ }
\newcommand{\Subsection}{\subsection}
\providecommand{\nobreakdash}{\nobreak\hbox{-}}
\setlength{\emergencystretch}{2em}
\multlinegap0pt
\usepackage{amssymb}
\usepackage{enumerate}
\usepackage{tikz}
\usepackage{cancel}
\usepackage{mathtools}
\newtheorem{theorem}{Theorem}
\title{Criteria for the presence of the maximal ideal in the set of associated primes}
\author{Mehrdad Nasernejad, Jonathan Toledo}
\date{Source: arXiv:2510.16566v2 (CC BY 4.0)}
\begin{document}
\maketitle

\noindent\emph{Source and licence.} Criteria for the presence of the maximal ideal in the set of associated primes. Version arXiv:2510.16566v2 (CC BY 4.0), distributed by arXiv under the Creative Commons Attribution 4.0 International licence, http://creativecommons.org/licenses/by/4.0/, as recorded in the arXiv licence metadata for this exact version and shown on the abstract page https://arxiv.org/abs/2510.16566v2. Full licence notice: \texttt{licenses/arxiv-2510.16566-CC-BY-4.0.txt}.\par\smallskip

\noindent\textit{Editorial scope.} Counted unit: the article's theorem Th.Not.maximal.1 (source0.tex lines 503--509). The article's complete original proof of it is delivered with it (source0.tex lines 511--530, one \texttt{proof} environment of 20 lines). No mathematics is written, repaired or re-derived: the statement and the proof are the article's own lines, in the article's own notation, and no retained line is substituted or re-flowed.  No further source-local context is needed: every cross-reference of every retained line resolves inside the two delivered spans.  Nothing was omitted from any retained span, so the package carries no omission record.  No retained line contains a citation, so the wrapper carries no bibliography and no bibliography entry.  No retained line contains a figure environment, an image-inclusion command or drawing code, so no image is delivered and the package declares no asset dependency.  Editorial formatting only, in addition: an AMS \texttt{amsart} preamble that restates the article's own declarations and macros used by the retained lines, namely the environments \texttt{theorem} on separate counters, together with the standard packages those lines need.  The retained environments print in this document as Theorem 1 in order of appearance, whereas the article prints the same items with the numbering of its own full document.  The complete native source of the article as distributed in the arXiv e-print for this exact version is bundled unchanged as \texttt{sources/187345/source0.tex}.
\begin{theorem} \label{Th.Not.maximal.1}
Suppose that    $I \subset R=K[x_1, \ldots, x_n]$ with $n\geq 2$,   is   a monomial ideal, $\mathfrak{m}=(x_1, \ldots, x_n)$, and 
  $\mathcal{G}(I)=\{x^{r_{1,1}}_1\cdots x^{r_{1,n}}_n, \ldots, x^{r_{k,1}}_1\cdots x^{r_{k,n}}_n\}$
  such that  there exist some  $1\leq i\neq j \leq n$ with 
  $$r_{1,i} \geq r_{2,i} \geq \cdots \geq r_{k, i} \quad \text{and} \quad r_{1,j} \geq r_{2,j} \geq \cdots \geq r_{k, j}.$$ 
   Then   $\mathfrak{m}\notin \mathrm{Ass}(R/I)$. 
\end{theorem}

\begin{proof} 
Without loss of generality, we may assume that   $r_{1,i} \geq r_{2,i} \geq \cdots \geq r_{k, i}$ for $i=1,2$. 
On the contrary, assume that $\mathfrak{m}\in \mathrm{Ass}(R/I)$. This implies that there exists some monomial $h\in R$ such that  $h\notin I$ and 
  $x_ih\in I$ for all $i=1, \ldots, n$. In particular, we have $x_1h, x_2h \in I$.   Let $h=\prod_{i=1}^n x_i^{\theta_i}$. We thus get 
 $x_1^{\theta_1 +1} \prod_{i=2}^n x_i^{\theta_i}\in I$ and $x_1^{\theta_1} x_2^{\theta_2+1}\prod_{i=3}^n x_i^{\theta_i} \in I$. 
  This gives  that there exist  $1\leq \lambda \leq k$  and $1\leq \gamma \leq k$  such that 
  \begin{equation} \label{15}
  x^{r_{\lambda,1}}_1\cdots x^{r_{\lambda,n}}_n \mid x_1^{\theta_1 +1} \prod_{i=2}^n x_i^{\theta_i} \;  \text {and} \;  
  x^{r_{\gamma,1}}_1\cdots x^{r_{\gamma,n}}_n \mid x_1^{\theta_1} x_2^{\theta_2+1}\prod_{i=3}^n x_i^{\theta_i}.
  \end{equation}
       By virtue  of    $\prod_{i=1}^n x_i^{\theta_i}\notin I$,   we must have
      \begin{equation} \label{16}
      x^{r_{\lambda,1}}_1\cdots x^{r_{\lambda,n}}_n \nmid  \prod_{i=1}^n x_i^{\theta_i} \;\; \text{and} \;\;  
     x^{r_{\gamma,1}}_1\cdots x^{r_{\gamma,n}}_n \nmid  \prod_{i=1}^n x_i^{\theta_i}.
      \end{equation}
     It follows from (\ref{15}) and (\ref{16})  that  $\theta_1 < r_{\lambda,1} \leq \theta_1+1$, $\theta_2 < r_{\gamma,2} \leq \theta_2+1$, 
   $r_{\lambda, 2} \leq \theta_2$, and $r_{\gamma, 1} \leq \theta_1$. Consequently, we  obtain $r_{\gamma, 1} \leq \theta_1 < r_{\lambda,1}$ and 
   $r_{\lambda, 2} \leq \theta_2< r_{\gamma,2}$. One can immediately deduce  from  our assumption  that 
    $\gamma < \lambda$ and $\lambda < \gamma$. This leads to a contradiction. Accordingly, we conclude that   $\mathfrak{m}\notin \mathrm{Ass}(R/I)$.    
\end{proof}

\end{document}
